\section{Introdução}
\label{introducao}

A evasão no ensino superior interrompe trajetórias pessoais, gera custo público evitável e reduz a eficiência das instituições \citep{realinho2022}. Desde \cite{tinto1975}, ela é entendida como resultado da interação entre o que o aluno traz ao ingressar e a sua integração acadêmica ao longo do curso, e não de uma causa única. O acompanhamento institucional, porém, costuma olhar indicadores isolados, como nota, idade ou situação financeira, que dizem quanto cada fator pesa, mas não como os fatores se combinam em cada aluno.

A unidade de análise deste trabalho é o estudante de graduação de uma instituição portuguesa, descrito no conjunto público UCI 697 \citep{realinho2022}. A trajetória é observada em três momentos em que a instituição pode agir: a matrícula, o fim do 1º semestre e o fim do 2º semestre. Os trabalhos sobre esse conjunto são tabulares, tratam cada aluno isoladamente e comparam classificadores \citep{realinho2022, villar2024}. Quando registros de alunos são convertidos em grafos de semelhança, o resultado depende muito de como o grafo é construído \citep{almeida2025}; falta comparar, sob a mesma avaliação, as duas decisões que definem esse grafo e o seu uso: como medir a semelhança entre alunos descritos por variáveis mistas e que tipo de método aplicar à estrutura resultante. Daí a questão de pesquisa: \textit{como a modelagem em grafos de características demográficas, socioeconômicas e acadêmicas apoia a identificação de perfis associados à evasão e ao sucesso, e qual combinação de distância e método faz isso melhor quando todas são avaliadas da mesma forma?}

Três hipóteses orientam a investigação. Pela homofilia \citep{mcpherson2001}, alunos semelhantes tendem a ter o mesmo desfecho, de modo que a estrutura do grafo deve carregar informação sobre a evasão acima do acaso (\textbf{H1}). Como a distância usual para dados mistos favorece certos tipos e escalas de variável \citep{vandevelden2026}, uma distância que corrige esse viés deve gerar um grafo que separa melhor os desfechos (\textbf{H2}). E, no mesmo grafo, um método que aprende com os atributos e a vizinhança de cada aluno deve prever o desfecho melhor que a atribuição pela comunidade, que troca resolução individual por interpretabilidade (\textbf{H3}).

A contribuição científica é comparar de forma controlada duas decisões de modelagem que costumam ser tomadas sem justificativa; a prática é oferecer perfis de risco interpretáveis, que permitem direcionar ações de permanência a grupos concretos, e indicar em que momento da trajetória o risco se torna visível. O objetivo geral é \textbf{desenvolver e avaliar um protótipo baseado em grafos que represente os alunos por semelhança e identifique perfis associados à evasão e ao sucesso}. Os objetivos teóricos são formalizar o grafo, comparar duas distâncias e dois métodos sob a mesma avaliação e medir como estrutura e desempenho mudam ao longo das janelas; os práticos, entregar um protótipo reprodutível e caracterizar perfis de risco e de sucesso.

A Seção~\ref{referencial} apresenta os conceitos e trabalhos relacionados; a Seção~\ref{desenvolvimento}, os dados, o grafo e a avaliação; a Seção~\ref{resultados}, os resultados e a sua discussão; e a Seção~\ref{conclusoes}, as conclusões. Código, figuras e relatório estão em \url{https://github.com/Erick-Bonifacio/student-dropout-success-graph}.
